\section{Conclusion}

We presented the \textbf{Confucius Code Agent (CCA)}, a coding agent for large-scale codebases. Built on the \textbf{Confucius SDK}, CCA separates and optimizes \emph{Agent Experience (AX)}, \emph{User Experience (UX)}, and \emph{Developer Experience (DX)}, enabling multi-step reasoning with modular tools, structured memory, and interpretable traces. Across public benchmarks and real-world settings, CCA shows that \emph{agentic scaffolding}—orchestration, memory, and tool abstractions—can matter as much as, or more than, the backbone model. The SDK’s hierarchical memory, context compression, and persistent notes support long-horizon stability, while extensions and the meta-agent facilitate rapid adaptation to new tools and workflows.

More broadly, Confucius SDK establishes a new foundation for AI agent research. Its modular architecture invites experimentation: from studying long-context reasoning and long-term memory, to exploring test-time adaptation, to integrating reinforcement learning with structured trajectory traces. We hope this framework accelerates progress toward AI developers that are powerful, interpretable, and continuously improving, bridging the persistent gap between research prototypes and the demands of real-world software engineering.
